\documentclass[12pt]{article}

\usepackage{amsmath, amssymb}
\usepackage{geometry}
\usepackage{graphicx}
\usepackage{hyperref}
\usepackage{xcolor}
\usepackage{booktabs}
\usepackage[numbers, sort&compress]{natbib}

\geometry{margin=1.2in}

\hypersetup{
    colorlinks=true,
    linkcolor=blue!60!black,
    citecolor=blue!60!black,
    urlcolor=blue!60!black
}

% Notation shortcuts
\newcommand{\Nzero}{N_0}
\newcommand{\sEdge}{s_e}
\newcommand{\sCore}{s_c}
\newcommand{\bEdge}{b_e}
\newcommand{\dEdge}{d_e}
\newcommand{\bCore}{b_c}
\newcommand{\dCore}{d_c}
\newcommand{\TODO}[1]{\textcolor{red}{\textbf{[TODO: #1]}}}

\title{Spatial rescue in a biofilm under antibiotic treatment:\\
a stochastic chain model}

\author{}
\date{}

\begin{document}
\maketitle

\begin{abstract}
Biofilms are spatially structured communities of bacteria in which cells near
the surface are exposed to antibiotic while interior cells are shielded. We
study how this structure affects the probability that a resistant mutant arises
and rescues the population before it goes extinct under lethal antibiotic
doses. We formulate a stochastic 1D ``chain model'' in which births and deaths
drive a conveyor-belt flow of cells between a protected core and an exposed
edge. We derive analytical approximations for the probability of mutant
appearance as a function of the initial population size, edge width, and
excess death rate, then validate and extend these results with stochastic
simulations. \TODO{Fill in main quantitative results once sweeps are complete.}
\end{abstract}

% --------------------------------------------------------------------------
\section{Introduction}
\label{sec:intro}
% --------------------------------------------------------------------------

Bacteria growing in biofilms are notoriously difficult to eradicate with
antibiotics~\cite{costerton1999}. One contributing factor is spatial structure:
cells embedded in a biofilm matrix are shielded from diffusing drug, so only the
outermost cells experience lethal concentrations. This creates an implicit refuge
that could, in principle, extend the window during which a resistant mutant can
arise and establish.

We ask a precise version of this question. Suppose a biofilm is exposed to an
antibiotic at a dose above the minimum inhibitory concentration (MIC), so that
surface cells die faster than they can replicate. The biofilm will eventually
collapse. Can the spatial structure give the population more time---and more
births---to produce a resistant mutant that rescues the lineage?

We compare two scenarios:
\begin{enumerate}
    \item \textbf{Well-mixed:} All $N_0$ cells experience the full antibiotic
        dose simultaneously.
    \item \textbf{Biofilm (chain model):} Only the outermost $l$ cells
        (the ``edge'') are exposed; the remaining $N_0 - l$ cells occupy a
        protected ``core.''
\end{enumerate}

The comparison is made under identical total population sizes and mutation
rates. Any difference in rescue probability is purely a consequence of spatial
structure.

% --------------------------------------------------------------------------
\section{Model}
\label{sec:model}
% --------------------------------------------------------------------------

\subsection{Population structure}

The biofilm is represented as a 1D array of $N(t)$ cells, indexed from $0$
(deep interior) to $N(t)-1$ (outermost surface). Two compartments are defined
by the parameter $l$ (edge width, in cells):
\begin{align}
    \text{Core} &= \text{cells } 0, 1, \ldots, b-1, \quad b \equiv \max(0,\, N-l), \\
    \text{Edge} &= \text{cells } b, b+1, \ldots, N-1.
\end{align}
When $N \leq l$ the core is empty and all cells are in the edge. The edge
always contains exactly $\min(N, l)$ cells; the core contains $\max(0, N-l)$.

\subsection{Birth--death dynamics (Gillespie algorithm)}

Events are drawn via the standard Gillespie algorithm (next-reaction
method). Each event type fires at a total rate equal to its per-capita rate
times the number of eligible cells. The per-capita rates are:

\begin{center}
\begin{tabular}{lllll}
\toprule
Compartment & Genotype & Birth rate & Death rate & Notes \\
\midrule
Edge  & WT & $\bEdge = 1$ & $\dEdge = \text{dose}$ & sweep variable \\
Edge  & R  & $b_R = 1$    & $d_R = 0.6$             & net positive growth \\
Core  & WT & $\bCore$     & $\dCore$                & $\bCore = \dCore$ at steady state \\
Core  & R  & $\bCore$     & $\dCore$                & same as core WT (neutral refuge) \\
\bottomrule
\end{tabular}
\end{center}

Time is measured in edge generations ($\bEdge \equiv 1$). The default
parameter values used in simulations are given in Table~\ref{tab:params}.

\textbf{Birth mechanics.} When any cell in compartment $c$ gives birth, a
daughter is inserted at one of the two positions immediately adjacent to the
parent, chosen with equal probability. Because the array shifts to accommodate
the new cell, the boundary $b = \max(0, N-l)$ increases by one: every birth
in the edge pushes the innermost edge cell (at position $b$) into the core,
and every birth in the core simply enlarges the core.

\textbf{Death mechanics.} When any cell dies it is removed from the array.
The boundary $b$ decreases by one: every death in the edge draws the outermost
core cell (at position $b-1$) into the edge, and every death in the core
simply shrinks the core.

\textbf{Conveyor-belt coupling.} The net effect is that cells flow inward
during edge births and outward during edge deaths. A core-born resistant
mutant can therefore reach the edge and eventually dominate it---the
\emph{core-delivery} rescue pathway---even without any directed migration.

\subsection{Mutation}

Resistance is modeled as an irreversible point mutation. With probability
$\mu$ per replication event, a daughter cell is assigned a new resistant (R)
lineage rather than inheriting the parent's genotype. Back-mutation is
ignored. Each de-novo lineage is tracked independently to determine its
origin compartment, time of appearance, and eventual fate.

\subsection{Rescue condition and simulation termination}

Treatment begins at $t = 0$ with the edge WT death rate jumping to
$\dEdge = \text{dose}$. Pre-treatment, $\dEdge = \bEdge = 1$ so the
population is at steady state. A simulation ends when one of the following
occurs:
\begin{itemize}
    \item \textbf{Rescue:} the number of R cells in the edge reaches the
        threshold $r^* = 10$.
    \item \textbf{Extinction:} $N = 0$.
    \item \textbf{Timeout:} a maximum number of events or elapsed time
        is reached.
\end{itemize}

\begin{table}[h]
\centering
\caption{Default simulation parameters.}
\label{tab:params}
\begin{tabular}{lll}
\toprule
Parameter & Value & Description \\
\midrule
$N_0$        & 1000 & Initial population size \\
$l$          & 100  & Edge width (biofilm); $l = 2000$ for well-mixed control \\
$\bEdge$     & 1.0  & WT edge birth rate (sets time unit) \\
$\bCore = \dCore$ & 0.2 & Core birth/death rates (neutral core) \\
$b_R$        & 1.0  & R edge birth rate \\
$d_R$        & 0.6  & R edge death rate ($s_R = 0.4$) \\
$\mu$        & $10^{-4}$ & Mutation probability per birth (main sweep) \\
$r^*$        & 10   & Rescue threshold (R edge cells) \\
\midrule
$\dEdge$     & 1.05--2.80 & Sweep variable (log-spaced in $s = \dEdge - 1$) \\
\bottomrule
\end{tabular}
\end{table}

% --------------------------------------------------------------------------
\section{Analytical approximations}
\label{sec:analytics}
% --------------------------------------------------------------------------

We derive approximate closed-form expressions for the probability that at
least one resistant mutant arises before the wild-type population goes extinct.
This is a necessary---but not sufficient---condition for rescue; the
additional requirements for establishment are analyzed in
Section~\ref{sec:turnover} and captured fully by the stochastic simulation.

\subsection{Well-mixed population}
\label{sec:wm_approx}

Consider first a population of $\Nzero$ cells all exposed to the same
antibiotic dose, with per-capita WT birth and death rates $\bEdge$ and
$\dEdge$. Define the net per-capita death rate:
\begin{equation}
    \sEdge \equiv \dEdge - \bEdge = \text{dose} - 1 > 0.
\end{equation}
In the deterministic limit the population declines exponentially:
\begin{equation}
    N(t) = \Nzero\, e^{-\sEdge t},
    \label{eq:wm_traj}
\end{equation}
reaching effective extinction ($N = 1$) at time $T = \ln \Nzero / \sEdge$.

Because mutations are replication-linked (probability $\mu$ per birth),
the total mutation rate at time $t$ is
\begin{equation}
    \lambda(t) = \mu\, \bEdge\, N(t).
\end{equation}
Mutations arrive as an inhomogeneous Poisson process, so the probability of
\emph{no} mutant arising by time $T$ is
\begin{equation}
    P(\text{no mutant}) = \exp\!\left(-\int_0^T \lambda(t)\, dt\right).
\end{equation}
Evaluating the integral:
\begin{equation}
    \int_0^T \mu\, \bEdge\, \Nzero\, e^{-\sEdge t}\, dt
    = \frac{\mu\, \bEdge\, \Nzero}{\sEdge}
      \left(1 - \frac{1}{\Nzero}\right)
    \approx \frac{\mu\, \bEdge\, \Nzero}{\sEdge},
\end{equation}
where the approximation holds for $\Nzero \gg 1$. Therefore:
\begin{equation}
    \boxed{P_{\text{WM}}(\text{mutant before extinction})
    = 1 - \exp\!\!\left(-\frac{\mu\, \bEdge\, \Nzero}{\sEdge}\right).}
    \label{eq:p_wm}
\end{equation}

The dimensionless group $\mu \bEdge \Nzero / \sEdge$ is the expected total
number of mutations produced over the entire decline. When it is much less
than one, $P_{\text{WM}} \approx \mu \bEdge \Nzero / \sEdge$; when much
greater than one, rescue is nearly certain.

\subsection{Biofilm: two-compartment, two-phase analysis}
\label{sec:biofilm_approx}

In the biofilm, the population dynamics differ from the well-mixed case
because core and edge cells have different rates and because the edge size
is fixed at $l$ cells as long as $N > l$.

\subsubsection*{Phase 1: core and edge coexist ($N > l$)}

While $N > l$, the core contains $N - l$ cells and the edge contains $l$
cells. The expected rate of change of total population size is:
\begin{equation}
    \frac{dN}{dt} = (\bCore - \dCore)(N - l) + (\bEdge - \dEdge)\, l
    = -\sCore (N - l) - \sEdge\, l,
    \label{eq:dNdt_phase1}
\end{equation}
where $\sCore \equiv \dCore - \bCore$. With the default parameters
$\bCore = \dCore = 0.2$, so $\sCore = 0$, and~\eqref{eq:dNdt_phase1}
simplifies to:
\begin{equation}
    \frac{dN}{dt} = -\sEdge\, l \quad (\sCore = 0),
    \label{eq:linear_decline}
\end{equation}
a \emph{linear} decline at a constant rate $v \equiv \sEdge\, l$ cells per
unit time. Intuitively, only the $l$ edge cells are dying faster than they
reproduce; the core maintains itself exactly and acts as a reservoir that
continuously replenishes the edge. The population reaches $N = l$ at time:
\begin{equation}
    T_1 = \frac{\Nzero - l}{\sEdge\, l}.
    \label{eq:T1}
\end{equation}

The total mutation rate during Phase~1 is:
\begin{equation}
    \lambda_1(t) = \mu\, [\bCore (N(t) - l) + \bEdge\, l]
    = \mu\, [\bCore (\Nzero - l - \sEdge\, l\, t) + \bEdge\, l].
\end{equation}
Integrating over Phase~1:
\begin{align}
    \int_0^{T_1} \lambda_1(t)\, dt
    &= \mu\left[\bCore \frac{(\Nzero - l)^2}{2\,\sEdge\, l}
       + \bEdge\, l\, T_1\right] \notag \\
    &= \mu\left[\frac{\bCore(\Nzero - l)^2}{2\,\sEdge\, l}
       + \frac{\bEdge(\Nzero - l)}{\sEdge}\right].
    \label{eq:Lambda1}
\end{align}

\subsubsection*{Phase 2: edge only ($N \leq l$)}

Once the core is exhausted the population follows exponential decline as in
the well-mixed case, but starting from $N = l$:
\begin{equation}
    N(t) = l\, e^{-\sEdge(t - T_1)}, \quad t \geq T_1.
\end{equation}
The mutation contribution from Phase~2 is, by direct analogy with
\eqref{eq:p_wm}:
\begin{equation}
    \int_{T_1}^{T_2} \lambda_2(t)\, dt \approx \frac{\mu\, \bEdge\, l}{\sEdge}.
    \label{eq:Lambda2}
\end{equation}

\subsubsection*{Combined probability}

Defining the total expected mutation count
$\Lambda \equiv \Lambda_1 + \Lambda_2$ (the sum of
\eqref{eq:Lambda1} and \eqref{eq:Lambda2}):
\begin{equation}
    \boxed{P_{\text{BF}}(\text{mutant before extinction})
    = 1 - e^{-\Lambda},}
    \label{eq:p_bf}
\end{equation}
\begin{equation}
    \Lambda = \mu\left[
        \frac{\bCore(\Nzero - l)^2}{2\,\sEdge\, l}
        + \frac{\bEdge(\Nzero - l)}{\sEdge}
        + \frac{\bEdge\, l}{\sEdge}
    \right]
    = \frac{\mu}{\sEdge}\left[
        \frac{\bCore(\Nzero - l)^2}{2\,l}
        + \bEdge\, \Nzero
    \right].
    \label{eq:Lambda_total}
\end{equation}

\subsubsection*{Comparison to well-mixed: effect of core birth rate on mutation supply}

The ratio $\Lambda / \Lambda_{\text{WM}}$ quantifies the biofilm's
mutation-supply advantage:
\begin{equation}
    \frac{\Lambda}{\Lambda_{\text{WM}}}
    = 1 + \frac{\bCore(\Nzero - l)^2}{2\, \bEdge\, \Nzero\, l}.
    \label{eq:ratio}
\end{equation}
This ratio is always $\geq 1$. For the default parameters ($\Nzero = 1000$,
$l = 100$, $\bCore = \bEdge = 1$) it equals $1 + 4.05 = 5.05$: the biofilm
expects five times as many mutations as the well-mixed control.

\paragraph{Dormant-core limit ($\bCore \to 0$).}
Setting $\bCore = 0$ in equation~\eqref{eq:Lambda_total}:
\begin{equation}
    \Lambda\big|_{\bCore=0}
    = \frac{\mu\,\bEdge\,(\Nzero - l)}{\sEdge}
      + \frac{\mu\,\bEdge\,l}{\sEdge}
    = \frac{\mu\,\bEdge\,\Nzero}{\sEdge}
    = \Lambda_{\text{WM}}.
    \label{eq:dormant_limit}
\end{equation}
A completely dormant core ($\bCore = \dCore = 0$) provides \emph{exactly as
many mutation opportunities as the well-mixed control.} The core extends the
population's lifetime---core cells are not killed by the antibiotic during
Phase~1---but without births it contributes zero mutations. The entire
mutation-supply advantage of the biofilm therefore comes from core births
alone; the surplus in equation~\eqref{eq:ratio} is exactly zero when
$\bCore = 0$ and grows linearly with $\bCore$ thereafter.

\subsection{Two effects of core metabolic activity on the core-delivery pathway}
\label{sec:turnover}

Core metabolic activity---the turnover rate $\rho \equiv \bCore = \dCore$
at which core cells divide and die---drives a fundamental tension in the
core-delivery rescue pathway. The mutation-supply analysis above captures
only one side of it:

\begin{description}
    \item[Effect 1 (mutation supply, captured by $\Lambda$).]
    Every core birth is a mutation opportunity. The core-delivery term in
    $\Lambda$ grows linearly with $\bCore$ (equation~\eqref{eq:ratio}),
    and vanishes in the dormant limit (equation~\eqref{eq:dormant_limit}).

    \item[Effect 2 (genetic drift, not captured by $\Lambda$).]
    A resistant lineage born in the core lives in a neutral environment
    ($b_{R,c} = \bCore$, $d_{R,c} = \dCore$): it has no growth advantage
    there. It must survive \emph{neutral genetic drift} long enough to be
    conveyed to the edge by the retreating boundary. Higher turnover means
    faster drift, so most lineages are lost before delivery.
\end{description}

We quantify Effect~2 and show how the two combine.

\subsubsection*{Lineage survival under neutral drift in the core}

A mutant born at depth $d$ from the core--edge boundary (i.e.\ $d$
positions away from being edge-exposed) must wait for the boundary to
retreat past it. Since the boundary retreats at rate $\sEdge\,l$ cells
per unit time (equation~\eqref{eq:linear_decline}), the expected waiting
time is
\begin{equation}
    \tau_d \approx \frac{d}{\sEdge\,l}.
    \label{eq:delivery_time}
\end{equation}
During this wait, the lineage undergoes a critical ($\bCore = \dCore$)
birth-death process. Starting from a single cell, the probability of
surviving to time $\tau$ is \cite{athreya1972}
\begin{equation}
    P(\text{survive to delivery} \mid d)
    = \frac{1}{1 + \bCore\,\tau_d}
    = \frac{1}{1 + \bCore\,d\,/\,(\sEdge\,l)}.
    \label{eq:survival_prob}
\end{equation}
Since births are placed at uniformly random depths $d \in [0,\,\Nzero - l]$
within the core, the mean survival probability over all birth positions is
\begin{equation}
    \langle P(\text{survive to delivery}) \rangle
    = \frac{\ln(1 + \alpha)}{\alpha}, \qquad
    \alpha \equiv \frac{\bCore\,(\Nzero - l)}{\sEdge\,l},
    \label{eq:mean_survival}
\end{equation}
where $\alpha$ is the dimensionless ratio of core turnover to boundary
retreat rate. In the two limits:
\begin{align}
    \langle P \rangle &\;\to\; 1
        && \alpha \to 0
        \;\;\text{(dormant: every lineage survives until delivery)},
        \label{eq:dormant_survival}\\[4pt]
    \langle P \rangle &\;\to\; \frac{\ln\alpha}{\alpha} \;\to\; 0
        && \alpha \to \infty
        \;\;\text{(very active: most lineages lost to drift)}.
        \label{eq:active_survival}
\end{align}

\subsubsection*{Net rate of core lineages reaching the edge}

Multiplying the total core mutation rate by the mean survival probability
gives the rate at which core mutant lineages are successfully delivered
to the edge:
\begin{align}
    \Gamma_{\text{core}}
    &= \underbrace{\mu\,\bCore\,(\Nzero - l)}_{\text{Effect 1: mutation rate}}
       \;\times\;
       \underbrace{\frac{\ln(1+\alpha)}{\alpha}}_{\text{Effect 2: mean survival}}
    \notag \\[6pt]
    &= \mu\,\sEdge\,l\;\ln\!\left(1 + \frac{\bCore(\Nzero-l)}{\sEdge\,l}\right).
    \label{eq:core_delivery_rate}
\end{align}
In the two limits:
\begin{align}
    \Gamma_{\text{core}} &\approx \mu\,\bCore\,(\Nzero - l)
        && \alpha \to 0
           \;\;\text{(linear in }\bCore\text{; drift negligible)},
        \label{eq:gamma_dormant}\\[4pt]
    \Gamma_{\text{core}} &\approx \mu\,\sEdge\,l\;\ln\!\!\left(
        \frac{\bCore(\Nzero-l)}{\sEdge\,l}\right)
        && \alpha \to \infty
           \;\;\text{(logarithmic saturation)}.
        \label{eq:gamma_active}
\end{align}
At high turnover, increasing $\bCore$ yields only logarithmic gains in
delivered lineages: the additional mutations from faster turnover are
nearly all eliminated by drift before reaching the edge. Effects~1 and~2
nearly cancel.

\subsubsection*{Bolus delivery and the rescue threshold}

A lineage that survives to delivery does not arrive at the edge as a
single cell. For a critical birth-death process conditioned on survival
to time $\tau$, the expected lineage size at delivery is \cite{athreya1972}
\begin{equation}
    \mathbb{E}[\,\text{size} \mid \text{survive}\;\tau\,] = 1 + \bCore\,\tau.
    \label{eq:bolus_size}
\end{equation}
Surviving lineages from an active core therefore arrive as a multi-cell
\emph{bolus}. The martingale property of critical birth-death processes---the
expected lineage size including extinct lineages (counted as zero) is
preserved at all times---implies that Effects~2 and~3 cancel exactly
\emph{in expectation}:
\begin{equation}
    P(\text{survive}) \;\times\; \mathbb{E}[\,\text{size} \mid \text{survive}\,]
    \;=\;
    \frac{1}{1+\bCore\tau}\times(1+\bCore\tau) \;=\; 1
    \quad \text{for all }\bCore\text{ and }\tau.
    \label{eq:martingale}
\end{equation}
If rescue probability were linear in the number of cells delivered, turnover
would have no net effect on the core-delivery pathway. The rescue threshold
$r^{*} = 10$ is a discrete, nonlinear criterion: $k$ cells arriving
simultaneously are far more likely to cross it than $k$ single-cell arrivals
separated in time. The martingale cancellation therefore does not hold at
the level of rescue events, and an active core retains a residual advantage
through bolus delivery that $\Lambda$ does not capture.

\subsection{The dose-dependent transition: nursery vs.\ drift trap}
\label{sec:transition}

A direct consequence of the dormant-core limit
(equation~\eqref{eq:dormant_limit}) reframes the entire dose-dependence.
The passive buffer role of the core---extending the population's
lifetime---provides \emph{no rescue advantage} over the well-mixed case:
the longer Phase~1 duration exactly compensates for having fewer cells in
the edge at any given time, leaving $\Lambda_{\text{edge}} = \Lambda_{\text{WM}}$
regardless of dose. It follows that

\begin{quote}
\textit{The entire rescue advantage of the biofilm over the well-mixed
case comes from core births.}
\end{quote}

Whether those births translate into rescue depends on $\alpha$.
The result is a clean two-regime picture:

\begin{center}
\renewcommand{\arraystretch}{1.6}
\begin{tabular}{lcc}
\toprule
 & Low selection ($\alpha \gg 1$) & High selection ($\alpha \ll 1$) \\
\midrule
Dormant core ($\bCore = 0$)
    & $P_{\text{BF}} \approx P_{\text{WM}}$
    & $P_{\text{BF}} \approx P_{\text{WM}}$ \\
Active core ($\bCore > 0$)
    & $P_{\text{BF}} \approx P_{\text{WM}}$
    & $P_{\text{BF}} \gg P_{\text{WM}}$ \\
\bottomrule
\end{tabular}
\end{center}

\noindent At low selection the active core is no better than a dormant
one: core mutations arise but are lost to drift before delivery
($\langle P(\text{survive})\rangle \to 0$), and rescue probability
collapses toward the well-mixed value. At high selection the active core
is decisive: the boundary retreats fast enough that drift has little time
to act, and the core functions as an efficient nursery. A dormant core,
by contrast, contributes nothing to rescue at any dose. \emph{Note: this
table captures drift effects on the core-delivery pathway only. A second
suppression mechanism---WT resupply into the edge---is analyzed
separately in Section~\ref{sec:resupply} and can drive
$P_{\rm BF} < P_{\rm WM}$ at low selection even in the dormant-core
case.}

\paragraph{The irony.}
The same antibiotic selection pressure that endangers the population is
the mechanism that makes the core valuable. High $\sEdge$ simultaneously
accelerates the collapse of the wild-type edge \emph{and} drives the
boundary retreat that delivers core-born mutants before drift can
eliminate them. The biofilm converts the killing force of the antibiotic
into a delivery mechanism for resistance.

\paragraph{The crossover.}
The transition between regimes occurs at $\alpha = 1$:
\begin{equation}
    \sEdge^{*} = \frac{\bCore(\Nzero - l)}{l}.
    \label{eq:crossover}
\end{equation}
For default parameters, $\sEdge^{*} = 0.2 \times 900 / 100 = 1.8$
(dose $= 2.8$). The main sweep spans $\sEdge \in [0.05,\, 2.0]$, which
straddles the crossover: the majority of simulated conditions fall in the
efficient-nursery regime ($\sEdge < \sEdge^{*}$), while the upper portion
of the sweep ($\sEdge \in [1.8,\, 2.0]$) probes the transition and the
drift-dominated side. The region near and above $\sEdge^{*}$ is therefore
the most mechanistically interesting part of the sweep.

\paragraph{Prediction for the $\bCore$ sensitivity sweep.}
Since the biofilm rescue advantage depends on $\bCore$ only through
core-delivery, and core-delivery is drift-suppressed at low selection,
the rescue curves for different values of $\bCore$ should \emph{fan out}
with increasing dose:
\begin{itemize}
    \item At low dose: all $\bCore$ values give
        $P_{\text{rescue}} \approx P_{\text{WM}}$ --- curves collapse
        onto the well-mixed baseline.
    \item At high dose: $P_{\text{rescue}}$ increases strongly with
        $\bCore$ --- curves diverge from one another and from the
        well-mixed baseline.
\end{itemize}
The dose at which the curves begin to separate is set by
equation~\eqref{eq:crossover}, which predicts a different separation
point for each $\bCore$ value. This gives a concrete, parameter-free
test of the mechanism against the simulation sweep.

\subsection{WT resupply and suppression of edge-born mutants}
\label{sec:resupply}

The mechanisms analyzed so far---mutation supply and core-delivery
drift---concern core-born lineages. The conveyor-belt dynamics impose a
second, distinct suppression mechanism on mutations arising directly in
the edge, acting through a \emph{wild-type resupply flux} from the core.

\subsubsection*{The mechanism}

During Phase~1 the edge compartment always contains exactly $l$ cells.
Every edge death is immediately compensated by the outermost core cell
moving into the edge. Because the core is wild-type dominated throughout
Phase~1, each edge death injects a fresh WT cell into the edge. The
total WT resupply rate is
\begin{equation}
    m_{\rm WT} \approx \dEdge \cdot l
    \quad \text{(WT cells per unit time during Phase~1)}.
    \label{eq:wt_resupply}
\end{equation}

In the well-mixed model there is no protected reservoir: dying WT cells
are simply lost, and the WT pool shrinks monotonically. A resistant
mutant in the WM population therefore faces progressively weaker
competition as the WT decline. In the biofilm the WT pool in the edge
is actively \emph{maintained} by resupply from the core throughout
Phase~1.

\subsubsection*{Effect on R establishment}

Let $k$ denote the number of R cells in the edge (held at $l$ cells
during Phase~1). The conveyor-belt mechanics---births push the innermost
edge cell to the core, deaths draw the outermost core cell to the
edge---yield effective transition rates for $k$:
\begin{align}
    \text{Rate}(k \to k+1) &= b_R\, k\, (l-k)/l,
        \label{eq:up_rate} \\
    \text{Rate}(k \to k-1) &= k\left[\bEdge\,(l-k)/l + d_R\right].
        \label{eq:down_rate}
\end{align}
The downward rate contains two contributions: WT births that push R
cells into the core (rate $\bEdge(l-k)k/l$), and R deaths replenished
by WT arriving from the core (rate $d_R k$). Both effects are sustained
at a constant level throughout Phase~1 because the WT pool in the edge
is maintained by resupply. This reduces the effective establishment
probability $p_{\rm est}^{\rm eff}$ below the value
$1 - d_R/b_R = 0.4$ that applies in a declining WM population.

The factors of $k/l$ and $(l-k)/l$ in
\eqref{eq:up_rate}--\eqref{eq:down_rate} assume \emph{uniform mixing}
within the edge: the cell at position $b$ (which births always push to
the core) is treated as a uniform sample from the edge, so it is R with
probability $k/l$. This approximation is most inaccurate for newly
delivered core-born R cells. Because edge deaths draw the outermost core
cell into position $b$ (the innermost edge slot), a freshly delivered R
lineage occupies exactly the position that the next birth event will push
back to the core. The initial displacement risk for a single delivered R
cell is therefore much higher than $1/l$, meaning
\eqref{eq:up_rate}--\eqref{eq:down_rate} underestimate establishment
difficulty and the true $p_{\rm est}^{\rm eff}$ is lower than the
first-passage solution to this chain would predict. The approximation
improves as $k$ grows and R cells spread through the edge via
replication.

An exact expression for $p_{\rm est}^{\rm eff}$ would require solving the
first-passage problem for the spatially explicit process, which does not
simplify readily. The qualitative conclusion is that WT resupply suppresses
edge-born R establishment, and this suppression is active for the full
Phase~1 duration $T_1 = (\Nzero - l)/v$.

\subsubsection*{Dose dependence}

The suppression window has duration $T_1 = (\Nzero - l)/v$, where
$v = \sEdge\,l$ is the boundary retreat velocity. At low selection
($v$ small, $T_1$ large), each edge-born R lineage faces sustained WT
dilution before the core is exhausted; suppression is strong. At high
selection ($v$ large, $T_1$ small), Phase~2 dominates quickly and WT
resupply has little time to act.

Importantly, WT resupply occurs even when $\bCore = 0$ (dormant core),
because core cells move to the edge upon edge deaths regardless of
whether they divide. The dormant-core equality
$\Lambda\big|_{\bCore=0} = \Lambda_{\rm WM}$ is a statement about
\emph{mutation supply} only; WT resupply suppresses establishment
probabilities through a channel not captured by $\Lambda$, and could
cause $P_{\rm rescue}^{\rm BF} \leq P_{\rm WM}$ even in the dormant
limit.

\subsubsection*{Combined suppression prediction}

Adding WT resupply to the drift-suppression picture gives a
\emph{non-monotone} dose dependence for the biofilm rescue advantage:

\begin{center}
\renewcommand{\arraystretch}{1.6}
\begin{tabular}{lcc}
\toprule
 & Low selection ($v$ small, $T_1$ large)
 & High selection ($v$ large, $T_1$ small) \\
\midrule
Dormant core ($\bCore = 0$)
    & $P_{\rm BF} \lesssim P_{\rm WM}$\;\;(WT resupply)
    & $P_{\rm BF} \approx P_{\rm WM}$ \\
Active core ($\bCore > 0$)
    & $P_{\rm BF} < P_{\rm WM}$\;\;(drift $+$ resupply)
    & $P_{\rm BF} \gg P_{\rm WM}$ \\
\bottomrule
\end{tabular}
\end{center}

\noindent The crossover dose at which the biofilm transitions from
suppression to enhancement is governed by the interplay of core drift
(equation~\eqref{eq:crossover}) and WT resupply
(equation~\eqref{eq:wt_resupply}). Whether this crossover is
empirically detectable, and at what dose it occurs, is a question for
the stochastic simulation.

\subsection{What the analytical approximation captures}
\label{sec:captures}

Table~\ref{tab:scope} summarizes which effects are included in the
approximation $P = 1 - e^{-\Lambda}$ and which are not.

\begin{table}[h]
\centering
\caption{Scope of the analytical approximation $P_{\text{BF}} = 1 - e^{-\Lambda}$.}
\label{tab:scope}
\renewcommand{\arraystretch}{1.4}
\begin{tabular}{p{6.2cm} c c}
\toprule
Effect & Direction & In $\Lambda$? \\
\midrule
Core births increase mutation supply           & $+$ & \checkmark \\
Core lifetime extends edge mutation window     & $+$ & \checkmark \\
Drift eliminates core lineages before delivery & $-$ & $\times$   \\
Bolus delivery aids threshold crossing         & $+$ & $\times$   \\
WT resupply suppresses edge-born establishment & $-$ & $\times$   \\
Establishment probability of edge-born mutants & --- & $\times$   \\
\bottomrule
\end{tabular}
\end{table}

$\Lambda$ counts every mutation event produced, treating each as equally
capable of rescue. The true rescue probability is therefore bounded above:
\begin{equation}
    P(\text{rescue}) \;\leq\; 1 - e^{-\Lambda}.
    \label{eq:upper_bound}
\end{equation}
The gap is set by establishment probabilities. For the \emph{edge-delivery}
pathway, a single R cell born in the edge has establishment probability
$1 - d_R/b_R = 1 - 0.6 = 0.4$ (exact result for a supercritical
birth-death process \cite{athreya1972}), though this is attenuated by WT
resupply throughout Phase~1 (Section~\ref{sec:resupply}). For the
\emph{core-delivery} pathway, establishment is further attenuated by the
drift-survival factor $\ln(1+\alpha)/\alpha$ and modified nonlinearly by
bolus delivery. Both pathways are tracked explicitly in the stochastic
simulation (Section~\ref{sec:sims}).

% --------------------------------------------------------------------------
\section{Simulation results}
\label{sec:sims}
% --------------------------------------------------------------------------

\TODO{This section will present results from the OSG parameter sweep:
main rescue-vs-dose curves for all three conditions (biofilm, chainWM,
wellMixed), pathway breakdown (edge-origin vs core-delivery rescue), and
sensitivity analyses over $\mu$ and $\bCore$.}

\subsection{Rescue probability vs.\ dose}
\label{sec:rescue_vs_dose}

\TODO{Figure: $P(\text{rescue})$ vs.\ $s = \text{dose} - 1$ for the three
conditions. Expected result: biofilm $>$ chainWM $\approx$ wellMixed at
every dose, with largest gap at intermediate doses.}

\subsection{Rescue pathway breakdown}
\label{sec:pathways}

\TODO{Figure: fraction of biofilm rescues attributed to edge-origin vs.\
core-delivery. Expected result: core-delivery fraction rises with dose,
consistent with the core becoming more important as edge mutations become
less likely to survive.}

\subsection{Sensitivity to core rate ($\bCore$)}
\label{sec:core_sensitivity}

\TODO{Figure: $P(\text{rescue})$ vs.\ dose for $\bCore / \bEdge \in
\{0.05, 0.1, 0.2\}$. Equation~\eqref{eq:ratio} predicts the rescue advantage
scales as $\bCore (\Nzero - l)^2 / (2 l \Nzero \bEdge)$, so slower core
growth should reduce the advantage.}

\subsection{Sensitivity to mutation rate ($\mu$)}
\label{sec:mu_sensitivity}

\TODO{Figure: $P(\text{rescue})$ vs.\ dose for $\mu \in \{10^{-6},
10^{-5}, 10^{-4}, 10^{-3}\}$. In the rare-mutation regime
($\Lambda \ll 1$), $P(\text{rescue}) \propto \mu$, so the curves
should shift vertically but preserve their shape.}

% --------------------------------------------------------------------------
\section{Discussion}
\label{sec:discussion}
% --------------------------------------------------------------------------

\TODO{Discuss the mechanism: the core acts as a long-lived nursery for
mutations that can be delivered to the edge by conveyor-belt flow, and
as a buffer that slows the overall population decline and keeps the
total birth rate high. Compare to Fruet et al.\ 2025 and other
biofilm-rescue literature.}

% --------------------------------------------------------------------------
\appendix
\section{Derivation of the linear-decline approximation}
\label{app:linear}
% --------------------------------------------------------------------------

Here we verify that the $\sCore = 0$ case yields exactly linear decline.
With $\bCore = \dCore$, the core birth and death rates cancel and
equation~\eqref{eq:dNdt_phase1} becomes:
\begin{equation}
    \frac{dN}{dt} = -\sEdge\, l = \text{const}, \quad N > l.
\end{equation}
Integrating from $t = 0$: $N(t) = \Nzero - \sEdge\, l\, t$.

The total birth rate during Phase~1 is:
\begin{align}
    B(t) &= \bCore (N(t) - l) + \bEdge\, l \\
         &= \bCore (\Nzero - l - \sEdge\, l\, t) + \bEdge\, l.
\end{align}
Integrating $\mu B(t)$ over $[0, T_1]$ with $T_1 = (\Nzero - l)/(\sEdge l)$:
\begin{align}
    \Lambda_1 &= \mu\int_0^{T_1} B(t)\, dt \\
    &= \mu\left[\bCore\int_0^{T_1}(\Nzero - l - \sEdge\, l\, t)\, dt
       + \bEdge\, l\, T_1\right] \\
    &= \mu\left[\bCore \cdot \frac{(\Nzero - l)}{2}\cdot T_1
       + \bEdge\, l\, T_1\right] \\
    &= \mu\, T_1\left[\frac{\bCore(\Nzero - l)}{2} + \bEdge\, l\right].
\end{align}
Substituting $T_1$ and combining with $\Lambda_2 = \mu \bEdge l / \sEdge$
recovers equation~\eqref{eq:Lambda_total}.

\bibliographystyle{unsrtnat}
\bibliography{refs}

\end{document}
